\chapter{Discussione}\label{chap:discussion}


\noindent Questo capitolo raccoglie la lettura dei risultati e non introduce grandezze nuove: il
suo contenuto è di interpretazione, e le misure a cui si riferisce sono quelle presentate nel
capitolo precedente. Vi sono ricostruiti gli esiti delle ipotesi e il quadro complessivo che ne
emerge, dichiarati i limiti dell'evidenza raccolta, e collocato il lavoro rispetto alla letteratura
--- la riduzione classica di un insieme di regole e la selezione guidata dai metadati ---
indicando quali verifiche estenderebbero i risultati e quali il materiale disponibile non ha
consentito.

\section{Esiti delle ipotesi}\label{sec:hypotheses-outcomes}

La \autoref{tab:hypotheses-outcomes} raccoglie gli esiti. Le ipotesi sono quelle enunciate nella
\autoref{tab:hypotheses} e nella \autoref{tab:exploratory}; nessuna è stata modificata dopo la
misura.

\begin{table}[htbp]
\centering
\caption{Esiti delle ipotesi. ``Confermata'' indica che la previsione è stata confermata dai dati;
``falsificata'' che i dati la contraddicono. Per H1 il valore riportato è quello ottenuto dopo la
correzione del difetto D5 (\autoref{sec:d5}): con il codice precedente la riduzione misurata era
dello \SI{0.1}{\percent}.}
\label{tab:hypotheses-outcomes}
\footnotesize
\begin{tabular}{@{}cp{3.6cm}p{4.0cm}p{3.4cm}@{}}
\toprule
\# & ipotesi & previsione & esito \\
\midrule
H1 & \str{score-argmin} con veti protetti & riduzione $\ge \SI{30}{\percent}$ & \textbf{falsificata} (\SI{0.3}{\percent}) \\
H6 & \str{veto-cover} su tutte le famiglie & output identico & \textbf{confermata} \num{113}/\num{113} \\
H7 & verifica empirica di \Vtwo{} & zero discrepanze & \textbf{confermata} \\
H8 & \str{precision-priority} da solo & guadagno di qualità & \textbf{confermata} ($+96$ corrette su \num{30} configurazioni) \\
H11 & \str{precision-priority} altrove & guadagno su Glass e Restaurant & \textbf{falsificata} \\
H14 & guadagno attribuibile alla precisione & il casuale non produce lo stesso guadagno & \textbf{falsificata} \\
H16 & riduzione e tempo proporzionali & tempi minori & \textbf{falsificata} (fino a \num{251}$\times$) \\
\midrule
H2  & \str{score-argmin} altrove & copertura non cala & confermata \\
H3  & copertura esatta e gratuita & nessuna perdita & falsificata \\
H4  & conversione \texttt{B}$\to$\texttt{D} & guadagno gratuito & falsificata \\
H5  & riordino delle righe del dataset & guadagno & falsificata \\
H9  & riordino $+$ potatura & guadagno maggiore & sub-additiva \\
H10 & regolarizzazione $\alpha = 3$ & guadagno $\ge \SI{20}{\percent}$ & falsificata \\
H12 & la regolarizzata generalizza meglio & guadagno maggiore & falsificata \\
H13 & riordino entro il cluster & leva di qualità & falsificata \\
H15 & ordinamento del budget migliorabile & guadagno & falsificata \\
\bottomrule
\end{tabular}
\end{table}

\paragraph{Le ipotesi confermate riguardano la teoria.} H6, H7 e H8 sono state confermate senza
riserve, e riguardano la garanzia e la leva di qualità presa da sola: la condizione \Vtwo{}
sopravvive alla verifica indipendente, la sua applicazione non altera il comportamento del sistema
su alcuna configurazione, e il riordino per priorità di precisione produce un guadagno quando è
impiegato da solo. È il nucleo dimostrativo del lavoro.

\paragraph{Le ipotesi falsificate riguardano le euristiche e le attribuzioni.} H1, H11, H14 e H16
sono state falsificate, e la loro falsificazione ha un valore che va oltre il caso singolo. H1
mostra che la strategia esatta lato generazione non è utilizzabile con la politica di tutela,
perché la classe garantita è troppo piccola per produrre la riduzione attesa. H11 mostra che il
guadagno non generalizza. H14 mostra che il segnale che si riteneva responsabile del guadagno non
lo spiega: il riordino casuale produce lo stesso effetto, e l'interpretazione che era stata
formulata in una fase precedente del lavoro è stata corretta. H16 mostra che la riduzione non
implica il risparmio di tempo.

\paragraph{La lettura d'insieme.} Il quadro che emerge è coerente e si può riassumere in quattro
punti. \emph{Primo}: esiste una potatura esatta, ed è piccola --- la condizione \Vtwo{} è corretta
e verificata, ma la riduzione che consente è limitata, e il limite è intrinseco al criterio
(\autoref{sec:maxextra}). \emph{Secondo}: esiste una potatura che migliora il risultato, ed è
condizionata --- \str{A5b} riduce dell'\SI{84.2}{\percent} e migliora il risultato su Bridges,
mentre sugli altri dataset danneggia, e la condizione di applicabilità è ora nota e misurabile.
\emph{Terzo}: il guadagno non proviene dal segnale che si riteneva responsabile, e la spiegazione
corretta è il prodotto di due fattori, la quota di imputazioni ottenute per via esatta e il numero
di imputazioni effettuate. \emph{Quarto}: il tempo non segue la riduzione, e una strategia può
rimuovere quasi tutte le regole peggiorando insieme il risultato e la durata.

Va infine osservato che i tre interventi studiati --- la potatura esatta, la potatura aggressiva e
la potatura conservativa --- non sono ordinabili. La prima non altera mai il risultato ma riduce
poco; la seconda produce il guadagno maggiore ma soltanto su una famiglia; la terza non peggiora
quasi mai ma rinuncia, nella maggior parte delle configurazioni, a qualunque riduzione. La scelta
fra le tre dipende dall'obiettivo: garanzia, guadagno, o assenza di rischio.

% =======================================================================================
\section{Limiti dell'evidenza}\label{sec:res-limits}

I limiti che seguono sono dichiarati perché delimitano la portata delle conclusioni, e non perché
ridimensionino i risultati: ciascuno indica anche la verifica che servirebbe a superarlo.

\paragraph{Il criterio diagnostico è post-hoc.} Il meccanismo a due fattori e il criterio che ne
discende sono stati ricavati osservando i risultati su tre famiglie, non previsti prima della
misura. La prova su una quarta famiglia è stata condotta (\autoref{sec:physician}) e ha prodotto
due esiti coerenti con il criterio, ma su \num{40} celle con copertura non nulla, a una sola
soglia, e senza che gli esiti siano distinguibili dal caso: due delle quattro versioni non sono
eseguibili per righe malformate nel dataset (\autoref{sec:d3}). Il criterio va quindi inteso come
una spiegazione con evidenza preliminare, non come una legge verificata su dati indipendenti.

\paragraph{Una sola famiglia con guadagno pieno.} Il guadagno di qualità è stato osservato su
Bridges; su Glass e Restaurant la stessa strategia danneggia, e la taratura conservativa riduce il
danno senza produrre guadagno. La generalità del meccanismo è sostenuta dalla coerenza dei tre
casi con la stessa legge, non da una molteplicità di famiglie favorevoli.

\paragraph{Le misure di tempo provengono da una sola macchina.} Le campagne sono state condotte su
una macchina a otto core, con esecuzioni strettamente sequenziali, e i valori assoluti dipendono
dall'ambiente. I confronti sono appaiati --- la stessa configurazione con e senza potatura --- e
sono quindi robusti rispetto all'ambiente, mentre i fattori di accelerazione e di rallentamento
riportati nel capitolo vanno letti come ordini di grandezza.

\paragraph{La numerosità è quella delle configurazioni.} L'unità di analisi è la configurazione,
non la cella nulla, e i test impiegano quindi alcune decine di osservazioni
(\autoref{sec:verification}). La potenza statistica è conseguentemente limitata: le conclusioni
sulle strategie il cui effetto è piccolo rispetto alla variabilità fra configurazioni vanno
considerate indicative. I test sono inoltre numerosi rispetto alle dimensioni del campione, e il
capitolo riporta per questo i valori corretti per confronti multipli accanto a quelli grezzi
(\autoref{sec:statistics}).

\paragraph{Il perimetro non comprende il confronto con altri metodi di imputazione.} Il lavoro
studia l'effetto della selezione delle regole \emph{a parità di sistema}: il termine di paragone è
la stessa esecuzione senza potatura, non un algoritmo di imputazione alternativo. La scelta
discende dal vincolo dichiarato --- la potatura opera sul file delle regole di un sistema
specifico --- e va tenuta presente nel leggere i risultati: essi non dicono se il sistema studiato
imputi meglio o peggio di altri, ma che cosa si guadagni o si perda potando le sue regole.

\paragraph{Due famiglie dell'archivio non sono utilizzabili.} La famiglia Nursery non contiene
dataset e la famiglia Physician, nelle versioni piene, ha dimensioni incompatibili con una
campagna (\num{103592} righe e fino a \num{121913} celle nulle). La griglia effettiva è quindi
composta da tre famiglie, ed è la ragione per cui la prova su una famiglia indipendente non ha
potuto che essere esplorativa.

% =======================================================================================
\section{Collocazione e verifiche ulteriori}\label{sec:positioning}

\paragraph{Rispetto alla riduzione classica di un insieme di regole.} La teoria classica delle
dipendenze funzionali affronta la ridondanza in forma dichiarativa: una dipendenza è superflua se
la sua rimozione non cambia l'insieme dei vincoli soddisfatti, e il calcolo di una copertura
minimale è un'operazione standard (\autoref{sec:rule-reduction}). Il lavoro mostra che quella
nozione non trasferisce al caso operativo, e la ragione è precisa: nel sistema studiato il ruolo di
una regola è definito da come il programma la consuma, non dalla sua semantica dichiarativa, e due
regole logicamente equivalenti possono avere effetti diversi perché differiscono per la direzione
delle soglie. La condizione \Vtwo{} è l'analogo operativo della dominanza classica, ma richiede
che la regola coprente sia \emph{più permissiva} a sinistra e \emph{più severa} a destra: la
direzione opposta a quella che si otterrebbe ragionando sui soli vincoli. Il teorema negativo
stabilisce che con l'informazione disponibile a una strategia esterna non si può andare oltre
questa classe senza rinunciare alla garanzia.

La distinzione non è soltanto argomentata: è misurata. Applicando il criterio classico come
strategia di potatura (\str{classical-cover}) sulle configurazioni di Bridges a soglia 9, la
riduzione raggiunge il \SI{97.4}{\percent} --- molto oltre il tetto della classe garantita --- ma
l'output cambia in tutte le configurazioni, l'accuratezza mediana perde \num{14.3} punti e le
imputazioni corrette scendono da \num{114} a \num{80} (\autoref{sec:res-classical}). Il confronto
mostra insieme due cose: che la ridondanza dichiarativa non implica la sicurezza operativa, e che
il numero di regole rimosse, da solo, non è un indicatore di qualità di una strategia.

\paragraph{Rispetto alla selezione guidata dai metadati.} Una linea di ricerca recente impiega
modelli linguistici per inferire relazioni fra colonne dal loro nome o per guidare
l'esplorazione di fonti eterogenee (\autoref{sec:semantic-pruning}). L'ambizione è la stessa ---
impiegare il significato invece della forma --- ma l'oggetto è diverso: quelle tecniche
\emph{interpretano} i metadati per produrne di nuovi, mentre il lavoro qui presentato
\emph{seleziona} metadati esistenti rispetto al consumo che ne fa un algoritmo noto. La differenza
ha una conseguenza pratica: la garanzia di equivalenza che il \autoref{chap:theory} dimostra non
dipende dalla qualità di una stima, ed è quindi verificabile per confronto di output; una
selezione guidata da un modello linguistico richiederebbe invece una validazione statistica del
modello stesso.

\uninatodo{Manca la misura dell'esercizio effettivo dei veti.}

\uninatodo{Manca il valore marginale del pruning esterno rispetto alla riduzione interna.}

\paragraph{Che cosa renderebbe più forti le conclusioni.} Tre verifiche sono state identificate
come decisive e non sono state condotte, ciascuna per una ragione dichiarata. La prima è la
\emph{misura dell'esercizio effettivo dei veti}: il meccanismo a due fattori è oggi inferito dalla
composizione degli output, e la strumentazione già disponibile consentirebbe di contare le
verifiche effettivamente eseguite prima e dopo la potatura, sostituendo l'inferenza con una misura
(\autoref{sec:future-work}). La seconda è il \emph{valore marginale della potatura esterna}
rispetto alla riduzione che il sistema compie già da sé sulle regole a soglia nulla con lato
sinistro chiave (\autoref{sec:dual-role}): quantificarlo richiede di calcolare la sovrapposizione
fra le due classi, e direbbe quanto del lavoro è già svolto dall'eseguibile. La terza è la
\emph{validazione prospettica del criterio} su una famiglia con quota esatta intermedia e
copertura stabile: la prova sulla quarta famiglia ha fornito due casi coerenti ma su \num{40}
celle, e non è sufficiente (\autoref{sec:physician}).

\paragraph{Che cosa non è stato possibile verificare.} Le configurazioni piene della famiglia
Physician restano fuori portata perché il preprocessing interno del sistema è quadratico nel numero
di righe, e le due versioni remodulate riparabili non sono state ripristinate perché il costo di
preprocessing non è giustificato dal numero di celle che se ne ricaverebbero
(\autoref{sec:d3}). Il confronto con metodi di imputazione alternativi è fuori dal perimetro
dichiarato del lavoro, che misura l'effetto della selezione delle regole a parità di sistema
(\autoref{sec:res-limits}). Entrambe le circostanze sono limiti del materiale o del perimetro, non
dell'impostazione, e sono state dichiarate prima che un lettore le incontrasse.
